% SOURCE_AUTHORITY: sga5_fr_workpass.tex, lines 14542--14573 (LF-normalized unit).
% SOURCE_SHA256: 791F4EFFC5E02832D5D77ED03518C8156D6F07E4C8238B03545DB93D883FBB28
\subsection*{§ 2. Correspondencia de Frobenius}

\subsubsection*{n\textsuperscript{o} 1. Definição}

Sejam $X$ um pré-esquema de característica $p$ e $F$ um feixe de conjuntos sobre $X_{\et}$. Para todo $X$-pré-esquema étale $U$ tem-se
\[
(\fr_X)_*(F)(U)=F(\fr_X^{-1}(U))=F(U^{(p/X)}),
\]
e o morfismo de Frobenius $\Fr_{U/X}:U\to U^{(p/X)}$ define uma aplicação
\[
F(\Fr_{U/X}):\fr_{X*}(F)(U)\longrightarrow F(U),
\]
manifestamente funtorial em $U$ (pois $\Fr_{U/X}$ é funtorial em $U$). Obtém-se assim um morfismo de feixes sobre $X_{\et}$:
\[
(\fr_X)_*(F)\longrightarrow F.
\]
Ora, quando $U$ é étale sobre $X$, $\Fr_{U/X}$ é um isomorfismo (§ 1 n$^\circ$2, prop. 2c)), o mesmo ocorre com $F(\Fr_{U/X})$; resulta que o morfismo precedente é um isomorfismo. Caso $F$ seja munido de uma estrutura algébrica (por exemplo, se for um feixe de $\Lambda$-módulos, sendo $\Lambda$ um anel de base dado), este isomorfismo é compatível com a estrutura algébrica. A inversa
\[
F(\Fr_{/X})^{-1}:F\longrightarrow \fr_{X*}(F)
\]
deste isomorfismo define, uma vez que $(\fr_X)^*$ é adjunto à esquerda do funtor $(\fr_X)_*$, um morfismo
\[
\Fr^*_{F/X}:(\fr_X)^*(F)\longrightarrow F.
\]
Pode-se dizer ainda mais, pois, sendo $\fr_X$ um homeomorfismo universal (§ 1 n$^\circ$2, prop. 2a)), os funtores $(\fr_X)^*$ e $(\fr_X)_*$ são quase-inversos um do outro e, por conseguinte, $\Fr^*_{F/X}$ é um isomorfismo, assim como o morfismo $F(\Fr_{/X})^{-1}$ do qual provém.

\begin{statement}{Definição 1}
Dado um pré-esquema $X$ de característica $p$ e um feixe de conjuntos $F$ sobre $X_{\et}$, chama-se correspondência de Frobenius sobre $(X,F)$ ao par $(\fr_X,\Fr^*_{F/X})$ formado pelo endomorfismo de Frobenius $\fr_X$ de $X$ e o isomorfismo
\[
\Fr^*_{F/X}=(F(\Fr_{/X})^{-1})^\sharp:(\fr_X)^*(F)\longrightarrow F.
\]
\end{statement}
